
\documentclass{article}
\usepackage{colortbl}
\usepackage{makecell}
\usepackage{multirow}
\usepackage{supertabular}

\begin{document}

\newcounter{utterance}

\centering \large Interaction Transcript for game `ifeval', experiment `ifeval', episode 82 with refineonly{-}llama{-}3.1{-}8b{-}instruct.
\vspace{24pt}

{ \footnotesize  \setcounter{utterance}{1}
\setlength{\tabcolsep}{0pt}
\begin{supertabular}{c@{$\;$}|p{.15\linewidth}@{}p{.15\linewidth}p{.15\linewidth}p{.15\linewidth}p{.15\linewidth}p{.15\linewidth}}
    \# & \multicolumn{2}{c}{Player} && \multicolumn{2}{c}{Game Master} \\
    \hline

    \theutterance \stepcounter{utterance}  
    & & & \multicolumn{4}{p{0.6\linewidth}}{
        \cellcolor[rgb]{0.9,0.9,0.9}{
            \makecell[{{p{\linewidth}}}]{
                \texttt{\tiny{[P1$\langle$GM]}}
                \texttt{A colt is 5 feet tall. It will grow 6 inches every month. How tall will it be in 3 years? Think step{-}by{-}step, then give your answer. Separate your thinking and the final answer by a line with just three "*" symbols: ***} \\
\texttt{At the end of your response, please explicitly add a postscript starting with P.P.S} \\
            }
        }
    }
     \\ \\

    \theutterance \stepcounter{utterance}  
    & & \multicolumn{4}{p{0.6\linewidth}}{
        \cellcolor[rgb]{0.95,0.95,0.95}{
            \makecell[{{p{\linewidth}}}]{
                \texttt{\tiny{[P1$\leftrightarrow$P]}}
                \texttt{{[}DRAFT{]}} \\
\texttt{To find the height of the colt in 3 years, we need to calculate the total growth in inches and then convert it to feet.} \\
\\ 
\texttt{1. First, we need to find out how many months are in 3 years. Since there are 12 months in a year, we multiply 3 years by 12 months/year.} \\
\texttt{   3 years * 12 months/year = 36 months} \\
\\ 
\texttt{2. Next, we need to find out how much the colt will grow in 36 months. We know that the colt grows 6 inches every month.} \\
\texttt{   36 months * 6 inches/month = 216 inches} \\
\\ 
\texttt{3. Now, we need to convert the total growth from inches to feet. There are 12 inches in a foot.} \\
\texttt{   216 inches / 12 inches/foot = 18 feet} \\
\\ 
\texttt{4. Finally, we need to add the total growth to the initial height of the colt.} \\
\texttt{   Initial height: 5 feet} \\
\texttt{   Total growth: 18 feet} \\
\texttt{   Total height: 5 feet + 18 feet = 23 feet} \\
\\ 
\texttt{***} \\
\texttt{The colt will be 23 feet tall in 3 years.} \\
\\ 
\texttt{P.P.S: To ensure accuracy, we should also consider that the colt's growth rate may not be constant, and there may be variations in growth depending on factors like nutrition, health, and genetics. However, based on the given information, the calculated height of 23 feet is an estimate.} \\
            }
        }
    }
    & \\ \\

    \theutterance \stepcounter{utterance}  
    & \multicolumn{4}{p{0.6\linewidth}}{
        \cellcolor[rgb]{0.561,0.737,0.561}{
            \makecell[{{p{\linewidth}}}]{
                \texttt{\tiny{[P1$\rangle$GM]}}
                \texttt{To find the height of the colt in 3 years, we need to calculate the total growth in inches and then convert it to feet.} \\
\\ 
\texttt{1. First, we need to find out how many months are in 3 years. Since there are 12 months in a year, we multiply 3 years by 12 months/year.} \\
\texttt{   3 years * 12 months/year = 36 months} \\
\\ 
\texttt{2. Next, we need to find out how much the colt will grow in 36 months. We know that the colt grows 6 inches every month.} \\
\texttt{   36 months * 6 inches/month = 216 inches} \\
\\ 
\texttt{3. Now, we need to convert the total growth from inches to feet. There are 12 inches in a foot.} \\
\texttt{   216 inches / 12 inches/foot = 18 feet} \\
\\ 
\texttt{4. Finally, we need to add the total growth to the initial height of the colt.} \\
\texttt{   Initial height: 5 feet} \\
\texttt{   Total growth: 18 feet} \\
\texttt{   Total height: 5 feet + 18 feet = 23 feet} \\
\\ 
\texttt{***} \\
\texttt{The colt will be 23 feet tall in 3 years.} \\
\\ 
\texttt{P.P.S: To ensure accuracy, we should also consider that the colt's growth rate may not be constant, and there may be variations in growth depending on factors like nutrition, health, and genetics. However, based on the given information, the calculated height of 23 feet is an estimate.} \\
            }
        }
    }
    & & \\ \\

    \theutterance \stepcounter{utterance}  
    & & & \multicolumn{2}{p{0.3\linewidth}}{
        \cellcolor[rgb]{0.9,0.9,0.9}{
            \makecell[{{p{\linewidth}}}]{
                \texttt{\tiny{[GM$|$GM]}}
                {'length_constraints:number_paragraphs': {'num_paragraphs': 2}, 'detectable_content:postscript': {'postscript_marker': 'P.P.S'}}
            }
        }
    }
    & & \\ \\

    \theutterance \stepcounter{utterance}  
    & & & \multicolumn{2}{p{0.3\linewidth}}{
        \cellcolor[rgb]{0.9,0.9,0.9}{
            \makecell[{{p{\linewidth}}}]{
                \texttt{\tiny{[GM$|$GM]}}
                \texttt{game\_result = WIN} \\
            }
        }
    }
    & & \\ \\

\end{supertabular}
}

\end{document}
